% Proposed unified replacement for the current Chapter 2 inputs.
% This file is intentionally not included by main.tex, so that it can be
% compared with the current chapter before any replacement is made.

\providecommand{\bW}{\mathbb{W}}

Throughout this chapter, let $C$ be a smooth, projective, geometrically connected
curve over a field $k$ of characteristic $p>0$. Fix a modulus
\[
    \mdls=\sum_i n_iP_i,
    \qquad U=C\setminus\mdls.
\]
For a single-point submodulus $\ndls=n_iP_i$, put
$d_\ndls=\deg(\ndls)$. The divisor $\ndls$ determines a $k$-rational point
$Z_\ndls\in C^{(d_\ndls)}$. We write
\[
    \pi_\ndls:X_\ndls:=\Bl_{Z_\ndls}\bigl(C^{(d_\ndls)}\bigr)
      \longrightarrow C^{(d_\ndls)}
\]
for the blowup, $E_\ndls$ for its exceptional divisor, and $\eta_\ndls$ for
the generic point of $E_\ndls$. By
\Cref{prop:relative_symmetric_power_smooth,prop:exceptional_divisor_smooth_point},
$E_\ndls\cong\bP_k^{d_\ndls-1}$ and in particular is an irreducible prime
divisor. Thus
\begin{equation}\label{equation:blowup_diagram_ndls}
\begin{tikzcd}
E_\ndls=\overline{\{\eta_\ndls\}}
  \arrow[r,hook]\arrow[d]&X_\ndls\arrow[d,"\pi_\ndls"]\\
Z_\ndls\arrow[r,hook]&C^{(d_\ndls)}.
\end{tikzcd}
\end{equation}

The purpose of the chapter is to prove
\Cref{theorem:SymmetricPowerOfSheavesIsTamelyRamifiedReduction}. If $\cF$ is
the rank-one $\Lambda$-local system in that theorem and
$\cP=\operatorname{Fr}(\cF)$ is its $G=\Lambda^\times$-torsor of frames, then
\Cref{prop:torsor_module_equivalence,definition:symmetric_power_local_system}
gives
\[
    \operatorname{Fr}\bigl(\cF^{[d_\ndls]}\bigr)
       \cong \cP^{[d_\ndls]}.
\]
The local system and its frame torsor therefore have the same ramification. It
is enough to prove the following torsor statement.

\begin{theorem}\label{theorem:torsors_after_blowup_is_tame}
Let $G$ be a finite abelian group and let $\cP\to U$ be a $G$-torsor whose
ramification is bounded by $\mdls$. For every single-point submodulus
$\ndls=n_iP_i$, the symmetric-power torsor $\cP^{[d_\ndls]}$ is tamely
ramified at $\eta_\ndls$.
\end{theorem}

The organization is deliberately parallel. We first define ramification for
finite covers and torsors. We then place Kummer, Artin--Schreier, and
Artin--Schreier--Witt theory next to one another, separating this general
background from the calculation particular to the thesis. A common prelude
collects the formal-local reduction, the contracted-product mechanism, and the
geometry of the exceptional valuation. Finally, the three cyclic calculations
are presented in the same order and with the same internal steps. The general
finite abelian case follows by primary decomposition.

% =====================================================================
\section{Ramification of Finite Covers and Torsors}
% =====================================================================

\subsection{Discrete valuation fields}

\begin{definition}[\stackstag{09E4}]
\label{definition:weakly_unramified_and_tamely_ramified_DVR}
An extension of discrete valuation rings is an injective local homomorphism
$A\subset B$ of discrete valuation rings. If $\pi_A$ and $\pi_B$ are
uniformizers, then $\pi_A=u\pi_B^e$ for a unit $u\in B$ and a uniquely
determined integer $e\geq1$. The integer $e$ is the \emph{ramification index}.
The extension is \emph{weakly unramified} if $e=1$.
\end{definition}

Let $A$ have fraction field $K$, let $L/K$ be finite separable, and let $B$ be
the integral closure of $A$ in $L$. For the maximal ideals $\mathfrak m_i$ of
$B$ above $\mathfrak m_A$, write $e_i$ and $f_i$ for the ramification indices
and residue degrees. Then
\[
    [L:K]=\sum_i e_if_i.
\]
When $A$ is henselian, there is only one prime above $\mathfrak m_A$.

\begin{definition}[\stackstag{09E9}]
\label{definition:tamely_ramified_dvrs}
With the notation above, $L/K$ is
\begin{enumerate}
    \item \emph{unramified} if every $e_i=1$ and every residue extension is
    separable;
    \item \emph{tamely ramified} if every residue extension is separable and,
    when the residue characteristic is $p>0$, every $e_i$ is prime to $p$;
    \item \emph{totally ramified} if there is one prime above
    $\mathfrak m_A$ and its residue extension is trivial.
\end{enumerate}
\end{definition}

Now suppose $A$ is complete with perfect residue field, $L/K$ is finite
Galois, and $G=\Gal(L/K)$. Let $B$ be the valuation ring of $L$ and $v_B$ its
valuation. The lower ramification groups are
\[
    G_i=\{\sigma\in G:v_B(\sigma(b)-b)\geq i+1
       \text{ for every }b\in B\},\qquad i\geq-1.
\]
Thus $G_{-1}=G$, $G_0$ is inertia, $L/K$ is unramified exactly when
$G_0=1$, and it is tame exactly when $G_1=1$. The Herbrand function
\[
    \phi_{L/K}(u)=\int_0^u\frac{dt}{[G_0:G_t]}
\]
reindexes the filtration by $G^v=G_{\phi_{L/K}^{-1}(v)}$. The reason for using
upper numbering in this thesis is its compatibility with quotients: if
$H\triangleleft G$, then
\[
    (G/H)^v=G^vH/H.
\]

\subsection{Ramification along a divisor}

Let $X$ be locally Noetherian, let $U\subset X$ be dense, and let
$Y\to U$ be finite etale. If $Z\subset X\setminus U$ is a prime divisor with
generic point $\xi$ and $\Oo_{X,\xi}$ is a discrete valuation ring, put
$K_\xi=\Frac(\Oo_{X,\xi})$. The pullback of $Y$ to $\Spec K_\xi$ is the
spectrum of a finite separable $K_\xi$-algebra $L_\xi$.

\begin{definition}\label{definition:tamely_ramified_in_codim_1}
The cover $Y\to U$ is \emph{unramified}, respectively \emph{tamely
ramified}, at $(Z,\xi)$ if every field factor of $L_\xi$ is unramified,
respectively tamely ramified, with respect to $\Oo_{X,\xi}$. It is unramified,
respectively tamely ramified, in codimension one if this holds at every prime
divisor in $X\setminus U$.
\end{definition}

Return to the curve $C$. For $P\in C\setminus U$, put
$\widehat K_P=\Frac(\widehat\Oo_{C,P})$. If $\cP\to U$ is a $G$-torsor,
then
\[
    \cP\big|_{\Spec\widehat K_P}\cong\coprod_{G/H}\Spec F,
\]
where $F/\widehat K_P$ is Galois and $H\subseteq G$ is the stabilizer of a
connected component.

\begin{definition}\label{definition:local_ramification_boundness}
The extension $F/\widehat K_P$ has \emph{ramification bounded by $r$} if
$H^r=1$. The torsor $\cP$ has ramification bounded by $r$ at $P$ if its
connected field factors have this property. It has \emph{ramification bounded
by $\mdls=\sum n_iP_i$} if its ramification at every $P_i$ is bounded by
$n_i$.
\end{definition}

Equivalently, the local torsor is classified by a continuous homomorphism
\[
   \rho_P:G_{\widehat K_P}\longrightarrow G,
\]
and the bound is the condition $\rho_P(G_{\widehat K_P}^{\,r})=1$. This
character formulation also shows that the definition is unchanged after
passing to the geometric points over $P$.

\subsection{Operations on ramified torsors}

\begin{lemma}\label{lemma:bounding_ramification_of_contracted_product}
Let $G$ be finite abelian and let $\cP_1,\cP_2$ be $G$-torsors on a dense open
$U\subset X$. If their ramification at a prime divisor is bounded by $r_1$
and $r_2$, then the ramification of
$\cP_1\otimes^G\cP_2$ is bounded by $\max(r_1,r_2)$.
\end{lemma}
\begin{proof}
Over the completed local field the two torsors correspond to characters
$\rho_1,\rho_2:G_K\to G$, while their contracted product corresponds to
$\rho_1+\rho_2$. Both characters vanish on
$G_K^{\max(r_1,r_2)}$, so their sum does as well.
\end{proof}

\begin{lemma}\label{lemma:descend_ramification_along_etale}
Let $f:X\to Y$ carry a prime divisor $(Z_X,\xi_X)$ to a prime divisor
$(Z_Y,\xi_Y)$ and be etale at $\xi_X$. If $\cP$ is a finite abelian torsor on
a dense open of $Y$, then $f^{-1}\cP$ has ramification bounded by $r$ at
$\xi_X$ if and only if $\cP$ has ramification bounded by $r$ at $\xi_Y$.
\end{lemma}
\begin{proof}
The induced extension of completed discrete valuation fields is unramified.
The upper-numbering filtration of the absolute Galois group is preserved under
unramified base change, so the two local characters kill the same $r$-th
ramification group.
\end{proof}

If $H\subseteq G$, the quotient $\cP/H$ corresponds locally to the composite
$G_K\to G\to G/H$. Consequently ramification bounds descend to quotient
torsors. Together with the preceding two lemmas, this is the functoriality used
in the final reduction from cyclic groups to arbitrary finite abelian groups.

% =====================================================================
\section{Three Parallel Theories of Cyclic Torsors}
% =====================================================================

This section contains general background, not the symmetric-power calculation
specific to the thesis. The three theories have the same formal shape:
\[
\begin{array}{c|c|c|c}
\text{theory}&H&\varphi:H\to H&\ker(\varphi)\\ \hline
\text{Kummer}&\bG_m&z\mapsto z^e&\mu_e\cong\Z/e\Z\\
\text{Artin--Schreier}&\bG_a&z\mapsto z^p-z&\F_p\cong\Z/p\Z\\
\text{Artin--Schreier--Witt}&\bW_r&F-\mathrm{id}&W_r(\F_p)\cong\Z/p^r\Z.
\end{array}
\]
In every row, pulling back $\varphi$ along a section $a$ produces a torsor
$\varphi(X)=a$, and changing $a$ by an element of $\varphi(H)$ does not change
its class.

\subsection{\texorpdfstring{Kummer theory: order prime to $p$}{Kummer theory: order prime to p}}

Let $e$ be prime to $p$ and suppose the field $K$ contains $\mu_e$. The Kummer
sequence
\[
    1\longrightarrow\mu_e\longrightarrow\bG_m
      \xrightarrow{(-)^e}\bG_m\longrightarrow1
\]
identifies
\[
    H^1(K,\mu_e)\cong K^\times/(K^\times)^e.
\]
The class of $a\in K^\times$ is represented by
$K[T]/(T^e-a)$, with $\zeta\in\mu_e$ acting by $T\mapsto\zeta T$.

\begin{theorem}[Ramification in Kummer extensions,
{\cite[Proposition~1.83]{koch1997algebraic}}]
\label{theorem:kummer_ramification}
Let $K$ be a discretely valued field containing $\mu_e$, with residue
characteristic prime to $e$, and let $L=K[T]/(T^e-a)$. Write $a=\pi^m u$ with
$u$ a unit. Then the ramification index is
\[
    e(L/K)=\frac{e}{\gcd(e,m)}.
\]
If $e\mid m$, the extension is unramified or trivial according to the residue
class of $u$; if $\gcd(e,m)=1$, it is totally and tamely ramified of degree
$e$. Conversely, every cyclic extension of degree prime to $p$ over a field
containing the required roots of unity is Kummer.
\end{theorem}

\subsection{\texorpdfstring{Artin--Schreier theory: order $p$}{Artin--Schreier theory: order p}}

For a field $K$ of characteristic $p$, the Artin--Schreier sequence
\[
  0\longrightarrow\F_p\longrightarrow\bG_a
    \xrightarrow{\wp}\bG_a\longrightarrow0,
  \qquad \wp(z)=z^p-z,
\]
identifies
\[
    H^1(K,\Z/p\Z)\cong K/\wp(K).
\]
The class of $a\in K$ is represented by $K[X]/(X^p-X-a)$, with
$g\in\F_p$ acting by $X\mapsto X+g$.

\begin{theorem}[Ramification in Artin--Schreier extensions,
{\cite{Thomas2005}}]\label{theorem:artin_schreier_ramification}
Let $K=k((x))$, with $k$ perfect of characteristic $p$, and let $a$ be a
reduced representative of a class in $K/\wp(K)$: either $a$ is regular, or
$v_x(a)=-m<0$ with $p\nmid m$. If the class is nonzero and regular, the
corresponding extension is unramified. If $v_x(a)=-m<0$, it is totally
ramified of degree $p$ and its unique upper ramification break is $m$.
Every cyclic extension of degree $p$ arises in this way.
\end{theorem}

\subsection{\texorpdfstring{Artin--Schreier--Witt theory: order $p^r$}{Artin--Schreier--Witt theory: order p-power}}
\phantomsection\label{subsection:witt_vector_prerequisites}

We recall only the Witt-vector structure used later. For $n\geq0$, the
$n$-th $p$-typical Witt polynomial is
\[
    w_n(X_0,\dots,X_n)=\sum_{j=0}^n p^jX_j^{p^{n-j}}.
\]
There are unique integral polynomials defining addition and multiplication on
$A^r$ for which the ghost map $(w_0,\dots,w_{r-1})$ is a ring map. The
resulting ring is denoted $W_r(A)$.

\begin{theorem}[Witt; {\cite[II, \S6]{serreLF}}]
\label{theorem:witt_ring_structure}
The functor $A\mapsto W_r(A)$ is a commutative ring functor. Witt addition is
given by universal polynomials $S_n$, and the $n$-th component has triangular
form
\begin{equation}\label{eq:witt_addition_shape}
   (\vec a\oplus\vec b)_n=a_n+b_n+C_n(a_0,\dots,a_{n-1};b_0,\dots,b_{n-1}).
\end{equation}
The ghost map is injective when $p$ is a non-zero-divisor and bijective when
$p$ is invertible.
\end{theorem}

\begin{lemma}[Ghost principle]\label{lemma:ghost_principle}
Two universal polynomial operations on Witt vectors are equal if their ghost
components agree over an integral polynomial ring.
\end{lemma}
\begin{proof}
Over that polynomial ring $p$ is a non-zero-divisor, so the ghost map is
injective. The resulting polynomial identity remains true after arbitrary
base change.
\end{proof}

\begin{prop}[Shape of Witt addition]\label{prop:witt_addition_shape}
The carry polynomial $C_n$ in \Cref{eq:witt_addition_shape} involves only
components of index strictly less than $n$. The analogous statement holds for
Witt subtraction $\ominus$.
\end{prop}
\begin{proof}
Solve recursively for $S_n$ in the identity
$w_n(S_0,\dots,S_n)=w_n(\vec X)+w_n(\vec Y)$. The terms containing $X_n,Y_n$
contribute $p^n(X_n+Y_n)$; all remaining terms involve lower components.
\end{proof}

For an $\F_p$-algebra $A$, Frobenius is componentwise,
$F(a_0,\dots,a_{r-1})=(a_0^p,\dots,a_{r-1}^p)$. Verschiebung and truncation are
\[
 V(a_0,\dots,a_{r-2})=(0,a_0,\dots,a_{r-2}),\qquad
 t(a_0,\dots,a_{r-1})=(a_0,\dots,a_{r-2}).
\]

\begin{definition}\label{definition:witt_operators}
We use $F$, $V$, and $t$ for these three operators on finite-length
$p$-typical Witt vectors.
\end{definition}

\begin{prop}\label{prop:witt_multiplication_by_p}
The map $t$ is a ring homomorphism with kernel $V^{r-1}A$; the map $V$ is
additive and commutes with $F$; and multiplication by $p$ is $VF=FV$. In
particular,
\begin{equation}\label{eq:witt_p_to_r_minus_1}
 p^{r-1}(a_0,\dots,a_{r-1})=(0,\dots,0,a_0^{p^{r-1}}).
\end{equation}
\end{prop}

\begin{example}\label{example:witt_of_Fp}
The additive group $W_r(\F_p)$ is canonically $\Z/p^r\Z$; the unit
$(1,0,\dots,0)$ is a generator.
\end{example}

\begin{definition}\label{definition:asw_operator}
The Artin--Schreier--Witt operator is the additive endomorphism
\[
    \wp=F-\mathrm{id}:W_r(A)\longrightarrow W_r(A).
\]
For $r=1$ this is $z\mapsto z^p-z$.
\end{definition}

\begin{lemma}[Witt bookkeeping]\label{lemma:witt_bookkeeping}
For an $\F_p$-algebra $A$:
\begin{enumerate}[label=(\roman*)]
 \item $t$ commutes with $F$ and $\wp$, and $\ker(t)=V^{r-1}A$;
 \item $V^{r-1}$ is additive and $\wp V^{r-1}=V^{r-1}\wp$;
 \item addition in the top component is carry-free:
 \[
   \vec a\oplus V^{r-1}(c)=(a_0,\dots,a_{r-2},a_{r-1}+c).
 \]
\end{enumerate}
\end{lemma}
\begin{proof}
The first two assertions follow from the preceding proposition. For the last,
compare ghost components: $V^{r-1}(c)$ has zero ghost components below level
$r-1$ and last ghost component $p^{r-1}c$.
\end{proof}

The Witt-vector group scheme $\bW_r$ has underlying scheme $\A^r_{\F_p}$ and
group law given by Witt addition. The operator $\wp$ gives the exact sequence
\begin{equation}\label{eq:asw_lang_sequence}
0\longrightarrow\Z/p^r\Z\cong W_r(\F_p)
 \longrightarrow\bW_r\xrightarrow{\wp}\bW_r\longrightarrow0.
\end{equation}

\begin{lemma}[ASW covers are etale torsors]\label{lemma:asw_etale_torsor}
Let $A$ be an $\F_p$-algebra and $\vec g\in W_r(A)$. Then
\[
 B=A[\vec X]/(\wp\vec X\ominus\vec g)
\]
is finite free of rank $p^r$ and etale over $A$. Translation
$\vec X\mapsto\vec X\oplus\underline m$ by
$m\in W_r(\F_p)\cong\Z/p^r\Z$ makes $\Spec B$ a $\Z/p^r\Z$-torsor.
\end{lemma}
\begin{proof}
By triangularity, the $n$-th equation is monic of degree $p$ in $X_n$ and
only involves lower variables otherwise. Hence the monomials
$\prod X_n^{a_n}$, $0\leq a_n<p$, form a basis. The Jacobian is triangular
with diagonal entries $-1$, so the algebra is etale. Its translations are the
kernel $W_r(\F_p)$ of $\wp$.
\end{proof}

\begin{theorem}[Artin--Schreier--Witt theory; {\cite[\S\S3--4]{Thomas2005}}]
\label{theorem:asw_theory}
For a field $K$ of characteristic $p$ there is a canonical isomorphism
\[
 W_r(K)/\wp W_r(K)\xrightarrow{\sim}
 H^1(K,\Z/p^r\Z)=\Hom_{\mathrm{cont}}(G_K,\Z/p^r\Z).
\]
The class of $\vec g$ corresponds to the torsor
$\wp\vec X=\vec g$. Its order is the order of the image of the associated
character; in particular, the torsor is connected of degree $p^r$ exactly
when $[\vec g]$ has order $p^r$. Every cyclic extension of degree dividing
$p^r$ arises this way.
\end{theorem}

\begin{cor}\label{cor:asw_class_order_and_first_component}
If $\vec g=(g_0,\dots,g_{r-1})$ has $g_0=0$, then $p^{r-1}[\vec g]=0$.
Consequently every representative of a class of order $p^r$ has nonzero
first component.
\end{cor}
\begin{proof}
Apply \Cref{eq:witt_p_to_r_minus_1}.
\end{proof}

\begin{rem*}\label{rem:asw_local_totally_ramified}
If $K=k((x))$ and $k$ is algebraically closed, every finite extension of $K$
is totally ramified. Hence a class of order $p^r$ defines a cyclic totally
ramified extension of degree $p^r$.
\end{rem*}

\begin{lemma}[Isobaricity of Witt addition]\label{lemma:witt_addition_isobaric}
Give the $j$-th component weight $p^j$. Every monomial in the $n$-th component
of a Witt sum has total weight $p^n$. The same is true for an iterated Witt
sum.
\end{lemma}
\begin{proof}
Each ghost polynomial $w_n$ is isobaric of weight $p^n$. The recursion defining
$S_n$ preserves that weight, and the assertion for iterated sums follows by
substitution.
\end{proof}

\begin{definition}\label{definition:reduced_witt_vector}
For $K=k((x))$, put $m(g)=\max(0,-v_x(g))$. A vector
$\vec g=(g_0,\dots,g_{r-1})\in W_r(K)$ is \emph{reduced} if every nonzero pole
order $m(g_j)$ is prime to $p$.
\end{definition}

\begin{theorem}[Schmid--Witt break formula]\label{theorem:schmid_witt_break_formula}
Let $K=k((x))$, with $k$ perfect, and let a reduced vector $\vec g$ represent
a cyclic totally ramified extension of degree $p^r$. Its largest upper
ramification break is
\[
    u=\max_{0\leq j<r}p^{r-1-j}m(g_j).
\]
\end{theorem}

For $r=1$ this is the Artin--Schreier break formula. The weighted maximum is
the exact feature that replaces the single pole-order bound in the degree-$p$
case.

% =====================================================================
\section{Prelude to the Explicit Ramification Calculations}
% =====================================================================

This section isolates everything shared by the three calculations. Its main
point is that Kummer multiplication, Artin--Schreier addition, and Witt
addition are instances of a single contracted-product construction.

\subsection{A common torsor equation}

Let $S$ be a scheme, let $H$ be a commutative $S$-group scheme, and let
$\varphi:H\to H$ be a finite etale surjective homomorphism with constant
kernel $G$. For a section $a:X\to H$ over an $S$-scheme $X$, write
\[
    \mathcal T_a:=X\times_{a,H,\varphi}H.
\]
It is a $G$-torsor on $X$. We write $\star$ for the group law of $H$; thus
$\star$ means multiplication for $H=\bG_m$, ordinary addition for
$H=\bG_a$, and Witt addition for $H=\bW_r$.

\begin{prop}[The common cyclic product calculation]
\label{prop:parallel_cyclic_product_torsor}
Let $a_i:X\to H$, $1\leq i\leq d$, be sections. There is a canonical
isomorphism of $G$-torsors
\[
    \mathcal T_{a_1}\otimes^G\cdots\otimes^G\mathcal T_{a_d}
       \cong \mathcal T_{a_1\star\cdots\star a_d}.
\]
On equation coordinates it is induced by
\[
    (X_1,\dots,X_d)\longmapsto Y:=X_1\star\cdots\star X_d,
    \qquad
    \varphi(Y)=a_1\star\cdots\star a_d.
\]
If a group $\Gamma$ permutes the indices, then $Y$ and the resulting datum are
$\Gamma$-invariant, so the isomorphism is compatible with descent to the
quotient by $\Gamma$.
\end{prop}
\begin{proof}
The product $\prod_i\mathcal T_{a_i}$ is a $G^d$-torsor. Contracting along the sum
map $G^d\to G$ gives its quotient by the kernel $\KG$. The displayed map is
$\KG$-invariant because the kernel translations cancel under $\star$; it is
equivariant for the residual $G$-action. Its target is a $G$-torsor because it
is the pullback of $\varphi$. Hence it is an isomorphism by
\Cref{prop:torsor_morphism_iso}. Commutativity and associativity of $\star$
give the permutation-equivariance.
\end{proof}

The three concrete instances are therefore
\[
\begin{array}{c|c|c}
G&Y&\text{new defining datum}\\ \hline
\Z/e\Z&\prod_iX_i&\prod_i a(x_i)\\
\Z/p\Z&\sum_iX_i&\sum_i a(x_i)\\
\Z/p^r\Z&\bigoplus_i\vec X_i&\bigoplus_i\vec a(x_i).
\end{array}
\]
This is the precise sense in which the three symmetric-power calculations are
the same calculation.

\subsection{Formal-local invariance and the exceptional valuation}

From now on assume $k$ is algebraically closed. This is the setting in which
the theorem is used later, and it makes every point of $C$ rational. Fix
$P\in C(k)$, put $\ndls=dP$, and choose a uniformizer $x$ at $P$. Let
$x_1,\dots,x_d$ be the induced parameters on $C^d$, and let
$e_1,\dots,e_d$ be their elementary symmetric polynomials. Set
\[
 V=\Spec k[[e_1,\dots,e_d]][e_d^{-1}],\qquad
 W=\Spec k[[x_1,\dots,x_d]][(x_1\cdots x_d)^{-1}].
\]
For each $i$, let $q_i:W\to\Spec k((x))$ be given by $x\mapsto x_i$.

\begin{lemma}[Formal-local invariance]\label{lemma:formal_local_invariance}
With the notation above:
\begin{enumerate}
 \item $\widehat\Oo_{C^{(d)},dP}\cong k[[e_1,\dots,e_d]]$, and the preimage of
 $U^{(d)}$ is $V$;
 \item if $\eta_{dP}$ is the generic point of the exceptional divisor, then
 \begin{equation}\label{eq:complete_field_at_generic_point}
   \widehat\Oo_{X_{dP},\eta_{dP}}\cong\kappa[[e_d]],\qquad
   \widehat K_{\eta_{dP}}\cong\kappa((e_d)),
   \quad
   \kappa=k(e_1/e_d,\dots,e_{d-1}/e_d);
 \end{equation}
 \item for every finite abelian $G$-torsor $\cP\to U$, with local restriction
 $\cP_x$ to $\Spec k((x))$, there is a canonical isomorphism over $V$
 \[
   \cP^{[d]}\big|_V\cong
   \Xi\backslash
   \bigl(q_1^{-1}\cP_x\times_W\cdots\times_Wq_d^{-1}\cP_x\bigr).
 \]
\end{enumerate}
Consequently the ramification of $\cP^{[d]}$ at $\eta_{dP}$ depends only on
$d$ and the isomorphism class of $\cP_x$.
\end{lemma}
\begin{proof}
Completion of $C^d$ at $(P,\dots,P)$ is
$k[[x_1,\dots,x_d]]$. Taking $S_d$-invariants commutes with the flat
completion and gives
\[
 k[[x_1,\dots,x_d]]^{S_d}=k[[e_1,\dots,e_d]].
\]
The boundary divisor consisting of cycles containing $P$ pulls back to
$V(x_1\cdots x_d)=V(e_d)$, proving (1).

Blowups commute with flat base change. On the blowup chart where the
$e_d$-direction is nonzero, write $e_i=(e_i/e_d)e_d$ for $i<d$. At the
generic point of the exceptional divisor the ratios $e_i/e_d$ are
algebraically independent residue parameters and $e_d$ is a uniformizer.
This gives (2).

Finally, the formation of the finite quotient
$\cP^{[d]}=\Xi\backslash\cP^d$ commutes with flat base change: the quotient
ring is an invariant subring, hence a kernel, and kernels commute with flat
base change. Pulling $\cP^d$ to $W$ gives the displayed product of the local
torsors. Restricting further along
$\Spec\kappa((e_d))\to V$ proves the consequence.
\end{proof}

\begin{cor}[Reduction to $\bP^1$]\label{cor:reduction_to_P1}
Suppose $\cP\to C\setminus\mdls$ and
$\cP'\to C'\setminus\mdls'$ are $G$-torsors and their restrictions to the
punctured formal discs at rational points $P,P'$ are isomorphic after matching
uniformizers. Then $\cP^{[d]}$ and $\cP'^{[d]}$ have identical ramification at
$\eta_{dP}$ and $\eta_{dP'}$.
\end{cor}
\begin{proof}
Both local symmetric-power torsors are obtained from the same local torsor by
the construction in \Cref{lemma:formal_local_invariance}.
\end{proof}

We will repeatedly use the following valuation interpretation. If a symmetric
Laurent polynomial belongs to
\[
   \kappa=k(e_1/e_d,\dots,e_{d-1}/e_d),
\]
then it is a unit or zero in $\kappa[[e_d]]$ and hence has no pole along the
exceptional divisor. By contrast, a power $e_d^a$ has valuation $a$.

\subsection{\texorpdfstring{Polynomial normal forms on $\bG_m$}{Polynomial normal forms on Gm}}

The formal-local reduction becomes useful once a local cyclic torsor is
realized by an equation on $\bG_m$.

\begin{lemma}[Surjectivity on integral Witt vectors]
\label{lemma:wp_surjective_on_integral_witt}
For algebraically closed $k$, the map
$\wp:W_r(k[[x]])\to W_r(k[[x]])$ is surjective.
\end{lemma}
\begin{proof}
For $r=1$, solve $u^p-u=h$ coefficientwise; algebraic closedness solves the
constant term. Induct on $r$. Solve after truncation, lift the solution to
length $r$, and use \Cref{lemma:witt_bookkeeping} to correct the remaining top
component by an element of $V^{r-1}k[[x]]$.
\end{proof}

\begin{definition}\label{definition:reduced_polynomial_vector}
A Witt vector $\vec f=(f_0,\dots,f_{r-1})\in W_r(k((x)))$ is a
\emph{reduced polynomial vector} if it is reduced and every
$f_j\in x^{-1}k[x^{-1}]$. Put $m_j=\deg_{x^{-1}}f_j$, with $m_j=0$ when
$f_j=0$.
\end{definition}

\begin{prop}[Reduced polynomial representatives]
\label{prop:reduced_polynomial_representatives}
Every class in $W_r(k((x)))/\wp W_r(k((x)))$ has a reduced polynomial
representative.
\end{prop}
\begin{proof}
Induct on $r$. For $r=1$, remove the integral part using surjectivity on
$k[[x]]$. If the remaining leading pole $cx^{-m}$ has $p\mid m$, subtract
\[
   \wp(c^{1/p}x^{-m/p})=cx^{-m}-c^{1/p}x^{-m/p};
\]
this strictly lowers the pole order, and iteration produces a reduced
polynomial.

For the induction step, first reduce the truncation to length $r-1$ and lift
the change of representative. The lower components are now reduced
polynomials. By \Cref{lemma:witt_bookkeeping}, subtracting
$\wp(V^{r-1}z)$ changes only the last component, by $-\wp(z)$. Apply the
length-one procedure to that component. Every operation subtracts an actual
element of $\wp W_r(k((x)))$, so the class is unchanged.
\end{proof}

\begin{lemma}[Parallel realization over $\bG_m$]
\label{lemma:realization_over_Gm}\label{lemma:realization_over_Gm_Wr}
Let $k$ be algebraically closed and let $\cQ$ be a connected cyclic torsor on
$\Spec k((x))$.
\begin{enumerate}
 \item If $G=\Z/e\Z$, $(e,p)=1$, then for some $a$ prime to $e$ the torsor is
 represented by
 \[
     T^e=x^a.
 \]
 The same equation defines a torsor on $\bG_m$.
 \item If $G=\Z/p\Z$ and the ramification is bounded by $d$, then the torsor is
 represented by
 \[
     X^p-X=f(x),\qquad f\in x^{-1}k[x^{-1}],\qquad
     m:=\deg_{x^{-1}}f<d,\qquad p\nmid m,
 \]
 with the convention $f=0$ for the trivial class. The same equation defines a
 torsor on $\bG_m$ extending etale across $\infty$.
 \item If $G=\Z/p^r\Z$ and the ramification is bounded by $d$, then the torsor
 is represented by
 \[
     \wp\vec X=\vec f,
 \]
 where $\vec f$ is a reduced polynomial vector satisfying
 \begin{equation}\label{eq:parallel_witt_pole_bound}
       p^{r-1-j}m_j<d\qquad(0\leq j<r).
 \end{equation}
 The same equation defines a torsor on $\bG_m$ extending etale across
 $\infty$.
\end{enumerate}
More generally, a tamely ramified $G$-torsor for finite abelian $G$ is induced
from the Kummer torsor in (1) along an injection $\Z/e\Z\hookrightarrow G$.
\end{lemma}
\begin{proof}
For (1), $k((x))$ has no nontrivial unramified extension and its tame quotient
is procyclic. Kummer theory identifies its connected quotient of order $e$
with the class of $x^a$ for $a\in(\Z/e\Z)^\times$.

For (2), Artin--Schreier theory identifies the class with an element of
$k((x))/\wp(k((x)))$. The length-one reduction in the proof of
\Cref{prop:reduced_polynomial_representatives} produces $f$ of the stated
form. The break formula gives $m<d$. Since $f$ is a polynomial in $x^{-1}$
without constant term, the equation is etale over $k[x^{-1}]$.

For (3), choose a reduced polynomial representative using
\Cref{prop:reduced_polynomial_representatives}. Connectedness and
\Cref{cor:asw_class_order_and_first_component} imply $f_0\neq0$. The
Schmid--Witt formula gives
\[
   \max_j p^{r-1-j}m_j<d,
\]
which is \Cref{eq:parallel_witt_pole_bound}. The torsor and its extension
across $\infty$ follow from \Cref{lemma:asw_etale_torsor} applied over
$k[x,x^{-1}]$ and $k[x^{-1}]$.
\end{proof}

\begin{rem*}\label{rem:realization_Wr_sanity}
The Witt pole bound forces $f_j=0$ whenever $p^{r-1-j}\geq d$. Thus the
strong restriction imposed by ``break $<d$'' is visible in every Witt
component, not only in the last one.
\end{rem*}

\subsection{Symmetric regularity}

Put $y_i=x_i^{-1}$. The identity
\[
   e_j(y_1,\dots,y_d)=\frac{e_{d-j}(x_1,\dots,x_d)}{e_d(x_1,\dots,x_d)}
\]
relates symmetric polynomials in the poles $y_i$ to the residue field
$\kappa$ of the exceptional valuation.

\begin{lemma}[Symmetric regularity]\label{lemma:symmetric_regularity}
If $\Phi\in k[y_1,\dots,y_d]^{S_d}$ and every monomial of $\Phi$ has total
degree strictly less than $d$, then
\[
    \Phi\in k(e_1/e_d,\dots,e_{d-1}/e_d)=\kappa.
\]
In particular, $\Phi$ is regular at $\eta_{dP}$.
\end{lemma}
\begin{proof}
Expand $\Phi$ in monomial symmetric functions. A symmetric function of degree
$<d$ is a polynomial in $e_1(y),\dots,e_{d-1}(y)$; $e_d(y)$ cannot occur.
The displayed identity puts each of these generators in $\kappa$.
\end{proof}

\subsection{The scheme-theoretic symmetric power causes no new divisor}

The torsor-theoretic symmetric power is a quotient of the ordinary symmetric
power of the covering scheme. The comparison is harmless in codimension one.

\begin{lemma}\label{lemma:torsor_unramified_codim_one}
Let $\cP=U'\to U$ be a finite abelian $G$-torsor on a smooth curve, and let
$C'$ be the smooth proper model of $U'$. The natural morphism from the ordinary
symmetric power $U'^{(d)}$ to the torsor-theoretic symmetric power
$U'^{[d]}$ is unramified along every codimension-one boundary component of
$C'^{(d)}$.
\end{lemma}
\begin{proof}
A boundary component has generic point represented in $C'^d$ by
$(Q,Q_2,\dots,Q_d)$, where the last $d-1$ entries are generic and independent.
The quotient group is $\Xi=\KG\rtimes S_d$. An element of the inertia group at
this generic point must fix each of the independent function-field factors,
so its permutation part and its last $d-1$ group components are trivial. The
condition that it lie in $\KG$ then forces its first component to be trivial.
Thus the inertia group is trivial, and the induced extension of discrete
valuation rings is weakly unramified.
\end{proof}

\section{Three Parallel Explicit Calculations}
% =====================================================================

We now perform the thesis-specific calculation. In all three cases we make the
same reductions and carry out the same five steps:
\[
\begin{array}{c|c|c|c}
&\Z/e\Z&(\Z/p\Z)&\Z/p^r\Z\\ \hline
\text{normal form}&T^e=x^a&X^p-X=f&\wp\vec X=\vec f\\
\text{product coordinate}&\prod_iT_i&\sum_iX_i&\bigoplus_i\vec X_i\\
\text{new datum}&e_d^a&\sum_i f(x_i)&\bigoplus_i\vec f(x_i)\\
\text{exceptional behavior}&v(e_d^a)=a&\text{regular}&\text{componentwise regular}\\
\text{conclusion}&\text{tame}&\text{unramified}&\text{unramified}.
\end{array}
\]

By \Cref{cor:reduction_to_P1,lemma:realization_over_Gm}, it is enough in each
case to take
\[
 C=\bP^1_k,\qquad U=\bG_m,\qquad \ndls=d\cdot0,
\]
and one of the displayed equations. The complete local field at the generic
point of the exceptional divisor is always
\[
   \widehat K_\eta=\kappa((e_d)),
   \qquad \widehat\Oo_\eta=\kappa[[e_d]],
   \qquad \kappa=k(e_1/e_d,\dots,e_{d-1}/e_d).
\]

\begin{theorem}\label{theorem:ramification_after_blowup_P1}
In the preceding model:
\begin{enumerate}
 \item if $\cP$ is a $\Z/e\Z$-torsor with $(e,p)=1$, then
 $\cP^{[d]}$ is tamely ramified at $\eta_{d\cdot0}$;
 \item if $\cP$ is a $\Z/p\Z$-torsor whose ramification at $0$ is bounded by
 $d$, then $\cP^{[d]}$ is unramified at $\eta_{d\cdot0}$;
 \item if $\cP$ is a $\Z/p^r\Z$-torsor whose ramification at $0$ is bounded by
 $d$, then $\cP^{[d]}$ is unramified at $\eta_{d\cdot0}$.
\end{enumerate}
\end{theorem}

\subsection{\texorpdfstring{The Kummer calculation: $G=\Z/e\Z$, $(e,p)=1$}{The Kummer calculation}}

\subsubsection*{Local normal form and realization}

By the realization lemma, the connected local torsor is represented by
\[
    T^e=x^a,
    \qquad \gcd(a,e)=1.
\]
The same equation over $k[x,x^{-1}]$ gives a $\Z/e\Z$-torsor on $\bG_m$.

\subsubsection*{Contracted product and descent}

Pull the equation back along the $d$ projections from $\bG_m^d$. By
\Cref{prop:parallel_cyclic_product_torsor}, the contracted product has
coordinate
\[
    Y=\prod_{i=1}^dT_i
\]
and equation
\[
    Y^e=\prod_{i=1}^d x_i^a=(x_1\cdots x_d)^a=e_d^a.
\]
Both $Y$ and $e_d^a$ are $S_d$-invariant. Hence this equation descends from
$\bG_m^d$ to the torsor $\cP^{[d]}$ on
\[
   \bG_m^{(d)}=\Spec k[e_1,\dots,e_d,e_d^{-1}].
\]

\subsubsection*{Restriction to the exceptional divisor}

Over $\widehat K_\eta=\kappa((e_d))$, the equation remains
\[
     Y^e=e_d^a.
\]
Its defining element has valuation $a$, and the ramification index divides
$e$. Since $(e,p)=1$, \Cref{theorem:kummer_ramification} shows that every
field factor is tamely ramified. Thus $\cP^{[d]}$ is tame at $\eta$.

\subsection{\texorpdfstring{The Artin--Schreier calculation: $G=\Z/p\Z$}{The Artin--Schreier calculation}}

\subsubsection*{Local normal form and realization}

The local torsor has an equation
\[
    X^p-X=f(x),
    \qquad f\in x^{-1}k[x^{-1}],
    \qquad m:=\deg_{x^{-1}}f<d.
\]
It extends to a torsor on $\bG_m$ that is unramified at $\infty$.

\subsubsection*{Contracted product and descent}

By the common product calculation, the $d$-fold contracted product on
$\bG_m^d$ has coordinate
\[
    Y=\sum_{i=1}^dX_i
\]
and equation
\[
    Y^p-Y=\alpha,
    \qquad \alpha:=\sum_{i=1}^df(x_i).
\]
The coordinate $Y$ and the datum $\alpha$ are $S_d$-invariant. Therefore
\[
  \cP^{[d]}\cong
  \Spec\frac{k[e_1,\dots,e_d,e_d^{-1}][Y]}{(Y^p-Y-\alpha)}.
\]

\subsubsection*{Degree estimate and exceptional restriction}

Writing $y_i=x_i^{-1}$, the expression $\alpha$ is a symmetric polynomial in
$y_1,\dots,y_d$ of total degree at most $m<d$. By
\Cref{lemma:symmetric_regularity},
\[
    \alpha\in\kappa\subset\widehat\Oo_\eta.
\]
Consequently the equation $Y^p-Y=\alpha$ defines a finite etale algebra over
$\widehat\Oo_\eta$: its derivative with respect to $Y$ is $-1$. Its generic
fiber is the restriction of $\cP^{[d]}$ to $\widehat K_\eta$, so every field
factor is unramified. This proves the degree-$p$ case.

\subsection{\texorpdfstring{The Artin--Schreier--Witt calculation: $G=\Z/p^r\Z$}{The Artin--Schreier--Witt calculation}}

\subsubsection*{Local normal form and realization}

First quotient away the unramified part of the local character. The remaining
inertia group is cyclic of order $p^s$ for some $s\leq r$; since the argument
treats all lengths at once, rename $s$ as $r$. We may therefore assume that the
local torsor is connected and totally ramified. By
\Cref{lemma:realization_over_Gm}, it has an equation
\[
     \wp\vec X=\vec f=(f_0,\dots,f_{r-1}),
     \qquad f_j\in x^{-1}k[x^{-1}],
\]
where, with $m_j=\deg_{x^{-1}}f_j$,
\[
     p^{r-1-j}m_j<d.
\]

\subsubsection*{Contracted product and descent}

Put $R=k[x,x^{-1}]$ and
\[
 S=R[\vec X]/(\wp\vec X\ominus\vec f).
\]
After pulling to $\bG_m^d$, define
\[
   \vec Y=\bigoplus_{i=1}^d\vec X_i,
   \qquad
   \vec\alpha=\bigoplus_{i=1}^d\vec f(x_i).
\]

\begin{prop}[Identification of the Witt product torsor]
\label{prop:witt_cover_identification}
Every component of $\vec Y$ is $\KG$-invariant, every component of
$\vec\alpha$ is a symmetric Laurent polynomial, and
\[
 (S^{\otimes_k d})^{\KG}
 \cong
 R^{\otimes_k d}[\vec Y]/(\wp\vec Y\ominus\vec\alpha).
\]
Taking $S_d$-invariants gives
\[
 \bigl((S^{\otimes_k d})^{\KG}\bigr)^{S_d}
 \cong
 k[e_1,\dots,e_d,e_d^{-1}][\vec Y]
      /(\wp\vec Y\ominus\vec\alpha).
\]
Thus
\[
  \cP^{[d]}\big|_{\Spec\widehat K_\eta}
   \cong
  \Spec\widehat K_\eta[\vec Y]/(\wp\vec Y\ominus\vec\alpha).
\]
\end{prop}
\begin{proof}
The $\KG$-translations telescope in the Witt sum $\vec Y$, so it is invariant,
and additivity of $\wp$ gives $\wp\vec Y=\vec\alpha$. The resulting morphism
from the displayed ASW torsor to the contracted product is equivariant between
$G$-torsors, hence an isomorphism by \Cref{prop:torsor_morphism_iso}; this is
also \Cref{prop:parallel_cyclic_product_torsor} for $H=\bW_r$.

Permuting the tensor factors preserves both iterated Witt sums because Witt
addition is commutative and associative. Since the action on Witt vectors is
componentwise, every component of $\vec Y$ and $\vec\alpha$ is invariant. The
ASW algebra is free over the base with basis
$\prod_nY_n^{a_n}$, $0\leq a_n<p$, and each basis element is invariant.
Therefore taking $S_d$-invariants acts only on the coefficients and produces
the second displayed algebra.
\end{proof}

\subsubsection*{Weighted degree estimate}

The Witt components are not ordinary componentwise sums, so the
Artin--Schreier degree estimate must be weighted.

\begin{theorem}[Regularity of all Witt components]
\label{theorem:witt_components_regularity}
Every component $\alpha_n$ of
$\vec\alpha=\bigoplus_i\vec f(x_i)$ belongs to $\kappa$. Consequently
$\vec\alpha\in W_r(\widehat\Oo_\eta)$.
\end{theorem}
\begin{proof}
Give $f_j(x_i)$ weight $p^j$. By
\Cref{lemma:witt_addition_isobaric}, every monomial
\[
   \prod_{i,j}f_j(x_i)^{a_{ij}}
\]
appearing in $\alpha_n$ satisfies
\[
   \sum_{i,j}a_{ij}p^j=p^n.
\]
Writing $A_j=\sum_i a_{ij}$ and expanding in $y_i=x_i^{-1}$, its total
$y$-degree is at most
\[
 \sum_jA_jm_j
   <d\sum_jA_jp^{j-(r-1)}
   =d\,p^{n-(r-1)}\leq d.
\]
The inequality is strict. The whole component $\alpha_n$ is symmetric by
\Cref{prop:witt_cover_identification}, so
\Cref{lemma:symmetric_regularity} gives $\alpha_n\in\kappa$.
\end{proof}

\subsubsection*{Restriction to the exceptional divisor}

By \Cref{theorem:witt_components_regularity}, the equation
\[
    \wp\vec Y=\vec\alpha
\]
is defined over $\widehat\Oo_\eta$. By
\Cref{lemma:asw_etale_torsor}, it is a finite etale
$\widehat\Oo_\eta$-algebra. Its generic fiber is precisely the restriction of
$\cP^{[d]}$ to $\widehat K_\eta$, hence every field factor is unramified.

\begin{proof}[Proof of \Cref{theorem:ramification_after_blowup_P1}]
The Kummer, Artin--Schreier, and Artin--Schreier--Witt subsections prove parts
(1), (2), and (3), respectively.
\end{proof}

\begin{cor}\label{cor:ramification_after_blowup_is_tame}
Let $\cP\to C\setminus\mdls$ be a finite abelian torsor and
$\ndls=dP$ a single-point submodulus.
\begin{enumerate}
 \item If $G=\Z/p\Z$ and the ramification at $P$ is bounded by $d$, then
 $\cP^{[d]}$ is unramified at $\eta_\ndls$.
 \item If $\cP$ is tamely ramified at $P$, then $\cP^{[d]}$ is tamely
 ramified at $\eta_\ndls$.
\end{enumerate}
\end{cor}
\begin{proof}
Use \Cref{lemma:realization_over_Gm} to match the punctured-formal-disc
restriction with the appropriate explicit torsor on $\bG_m$, then use
\Cref{cor:reduction_to_P1} and the first two calculations above. For a tame
$G$-torsor, induce the cyclic Kummer calculation along the injection of its
cyclic inertia group into $G$.
\end{proof}

\begin{theorem}\label{theorem:ramification_after_blowup_is_tame_for_p_power}
Let $G=\Z/p^r\Z$ and suppose the ramification of $\cP$ at $P$ is bounded by
$d=\deg\ndls$. Then $\cP^{[d]}$ is unramified at $\eta_\ndls$.
\end{theorem}
\begin{proof}
Quotient away the unramified part, apply the length-$r$ realization in
\Cref{lemma:realization_over_Gm}, transfer the calculation by
\Cref{cor:reduction_to_P1}, and use the Artin--Schreier--Witt calculation.
\end{proof}

% =====================================================================
\section{Assembly for a Finite Abelian Group}
% =====================================================================

We can now prove the torsor formulation stated at the beginning of the
chapter.

\begin{proof}[Proof of \Cref{theorem:torsors_after_blowup_is_tame}]
Fix $\ndls=dP$. Choose a decomposition
\[
   G\cong G_{p'}\times G_p,
   \qquad
   G_{p'}\cong\prod_i\Z/e_i\Z\quad((e_i,p)=1),
   \qquad
   G_p\cong\prod_j\Z/p^{r_j}\Z.
\]
Projecting the character of $\cP$ to these cyclic factors gives cyclic torsors
$\cP_i$ and $\cQ_j$. Conversely, $\cP$ is their product after extending
structure groups along the inclusions of the factors. Symmetric powers commute
with these products because both are obtained by contracting the product
torsor along the relevant homomorphism of finite abelian groups.

Every $\cP_i$ is tamely ramified at $P$, so its $d$-th symmetric power is tame
at $\eta_\ndls$ by the Kummer calculation. The ramification bound on $\cP$
descends to every quotient character; hence each $\cQ_j$ has ramification
bounded by $d$, and \Cref{theorem:ramification_after_blowup_is_tame_for_p_power}
shows that $\cQ_j^{[d]}$ is unramified at $\eta_\ndls$. Finally,
\Cref{lemma:bounding_ramification_of_contracted_product} shows that the product
of the tame prime-to-$p$ factors and the unramified $p$-primary factors is
tame. This product is $\cP^{[d]}$.
\end{proof}

% =====================================================================
\section{Ramification on Products of Blowups}
\label{subsection:ram_prod_analysis}
% =====================================================================

We finish with the compatibility needed in
\Cref{section:gcft}. Let $X$ and $Y$ be smooth $k$-schemes, let $x\in X$ and
$y\in Y$ be closed points, and write
\[
 \widetilde X=\Bl_x(X),\qquad
 \widetilde Y=\Bl_y(Y),\qquad
 \widetilde{X\times Y}=\Bl_{(x,y)}(X\times_kY).
\]
Denote the generic points of the three exceptional divisors by
$\eta_X,\eta_Y$, and $\eta_{X\times Y}$.

\begin{prop}\label{prop:ramification-external-product}
Let $G$ be a finite abelian group. Let $\cG_X$ and $\cG_Y$ be $G$-torsors on
open subsets of $\widetilde X$ and $\widetilde Y$ disjoint from the exceptional
divisors. If their ramification at $\eta_X$ and $\eta_Y$ is bounded by $r$,
then the external contracted product
\[
    \operatorname{pr}_1^{-1}\cG_X
       \otimes^G
    \operatorname{pr}_2^{-1}\cG_Y
\]
has ramification bounded by $r$ at $\eta_{X\times Y}$.
\end{prop}
\begin{proof}
The assertion is etale-local at $x$ and $y$, so by smoothness it is enough to
take $X=\A_k^n$, $Y=\A_k^m$, and $x=y=0$. Write the affine coordinates as
$x_1,\dots,x_n,y_1,\dots,y_m$. At the generic point of the exceptional divisor
of $\Bl_0(\A_k^{n+m})$, every coordinate is a uniformizer times a unit. On the
chart in which the $x_1$-direction is nonzero, the completed local ring is
\[
 k\left(\frac{x_2}{x_1},\dots,\frac{x_n}{x_1},
          \frac{y_1}{x_1},\dots,
          \frac{y_m}{x_1}\right)[[x_1]].
\]

There is a dense open neighborhood of $\eta_{X\times Y}$ mapping to both
$\widetilde X$ and $\widetilde Y$: it is the locus where neither the
$X$-directions nor the $Y$-directions vanish. The induced homomorphism from the
exceptional DVR of $\widetilde X$ to the displayed DVR sends the uniformizer
$x_1$ to a uniformizer and enlarges the residue field by purely transcendental
elements. It is therefore weakly unramified with separable residue extension.
The corresponding statement holds for $\widetilde Y$ (using any $y_j$ as
uniformizer on the overlap).

Upper ramification groups are preserved by such a weakly unramified,
residually separable extension. Hence both pulled-back torsors have
ramification bounded by $r$ at $\eta_{X\times Y}$. Their contracted product
has the same bound by
\Cref{lemma:bounding_ramification_of_contracted_product}.
\end{proof}
